Direct Preference Optimization (DPO) penalizes rejected responses as whole trajectories. This is a poor inductive bias for multi-step reasoning: a rejected solution can contain a correct prefix, a localized first error, and later steps that merely inherit the wrong state. We study causal rejected-side masking, a family of DPO-style objectives that changes only how the rejected response is aggregated. The central operation is prefix protection: steps before the first error are masked out of the rejected-side penalty. Causal-Masked DPO (\cmdpo) further adds a decayed penalty on post-error suffix steps. Controlled arithmetic experiments show that vanilla DPO strongly degrades reasoning accuracy, while causal masking improves retention and held-out performance. Across three harder-arithmetic seeds, \cmdpo{} is best on average, but only narrowly over a first-error-only baseline. Additional diagnostics show that truncating rejected traces suppresses errors but can still hurt correct prefixes, uniform down-weighting does not reproduce the same first-error suppression, and noisy first-error labels can break the intended mask. GSM8K-style pilots are mixed: prefix protection helps, but suffix decay has not yet produced a reliable additional gain. We therefore present \cmdpo{} as a causal masking framework, with prefix protection as the robust empirical finding and suffix decay as a promising but localization-sensitive design choice.
